\chapter{CLTV y métricas de cliente}
\label{ch:cltv}

\section{Fórmula del CLTV}

El enunciado fija el CLTV (\emph{Customer Lifetime Value}) por cliente
mediante la siguiente expresión, válida sobre una ventana temporal $t$:

\begin{align}
  \mathrm{CLTV}(c)
  \;=\;
  \text{Ingresos}_t(c) \;\times\;
  \text{Margen}_t(c) \;\times\;
  \text{Frecuencia}_t(c) \;\times\;
  R_t(c).
  \label{eq:cltv}
\end{align}

En la implementación del proyecto, la ventana es la cobertura efectiva
del dataset, $\,2020\text{-}01 \to 2026\text{-}01$, esto es, $W=72$ meses
naturales.

\section{Operacionalización}

Cada componente de la ecuación~\eqref{eq:cltv} se traduce a una columna
calculada por cliente en
\texttt{data/processed/cltv.parquet}. La
tabla~\ref{tab:cltv-componentes} formaliza esa traducción.

\begin{table}[H]
\centering
\caption{Operacionalización de la fórmula de CLTV sobre el DWH.}
\label{tab:cltv-componentes}
\small
\begin{tabular}{l p{9.5cm}}
\toprule
\textbf{Componente} & \textbf{Definición} \\
\midrule
$\text{Ingresos}_t$  &
  $\displaystyle
    \sum_{l \in \text{ventas}(c)} \text{subtotal}_l
    \;-\;
    \sum_{r \in \text{devoluciones}(c)} \text{returned\_value}_r$
  \;(columna \texttt{net\_revenue}). \\
$\text{Margen}_t$    &
  $\dfrac{\text{Ingresos}_t - \sum_l \text{unit\_cost}_l \cdot
    \text{quantity}_l}{\text{Ingresos}_t}$
  \;(columna \texttt{margin\_rate}, proporción $\in [-\,\dots\,,1]$). \\
$\text{Frecuencia}_t$ &
  número de tickets distintos:
  $|\{\,\text{sale\_id} : c\,\}|$
  \;(columna \texttt{frequency}). \\
$R_t$                 &
  retención: meses con al menos una compra dividido por la ventana,
  $R_t = \texttt{active\_months}/W$ con $W=72$
  \;(columna \texttt{retention\_rate}). \\
\bottomrule
\end{tabular}
\end{table}

\noindent\textbf{Comentarios.} El margen es una \emph{proporción} sobre los
ingresos netos, no un valor monetario; multiplicado por $\text{Ingresos}_t$
recupera la magnitud monetaria. Los clientes con $\text{Ingresos}_t < 0$
(devoluciones que superan ventas) o con coste mayor que ingresos pueden
tener $\text{Margen}_t$ negativo o $> 1$, lo cual se preserva tal cual sin
truncar para no falsear el dato original (véase
capítulo~\ref{ch:conclusiones}, \emph{Limitaciones}).

\section{Implementación}

La operacionalización anterior se reparte en dos lugares: una consulta SQL
en el notebook~03 que agrega los componentes por cliente, y el módulo
\texttt{src/analytics/cltv.py} que aplica la fórmula multiplicativa de la
ecuación~\eqref{eq:cltv}. El módulo entero es deliberadamente compacto:

\begin{lstlisting}[language=Python,
  caption={\texttt{src/analytics/cltv.py}: aplicación de la fórmula de
  CLTV sobre las componentes ya agregadas.},
  label={lst:cltv_module}]
"""Calculo de CLTV por cliente.

Formula del enunciado:
    CLTV = Ingresos_t x Margen_t x Frecuencia_t x R_t

Donde, en la ventana temporal t:
    - Ingresos_t   = ingresos netos del cliente (ventas - devoluciones).
    - Margen_t     = (ingresos - coste) / ingresos  (proporcion).
    - Frecuencia_t = n. de tickets distintos del cliente.
    - R_t          = retencion = meses_activos / meses_de_la_ventana.
"""
from __future__ import annotations
import pandas as pd

def compute_cltv(df : pd.DataFrame) -> pd.DataFrame:
    """Calcula el CLTV de cada cliente con compras."""
    df["cltv"] = (
        df["net_revenue"].astype(float)
        * df["margin_rate"].astype(float)
        * df["frequency"].astype(float)
        * df["retention_rate"].astype(float)
    )
    return df.sort_values("cltv", ascending=False).reset_index(drop=True)
\end{lstlisting}

\noindent El test
\texttt{test\_cltv\_formula\_consistency} (apéndice~\ref{ap:tests})
recalcula la fórmula sobre las primeras $50$ filas y comprueba que el
valor almacenado coincide con el producto de las cuatro componentes.

\section{Resultados agregados}
\label{sec:cltv-resultados}

La tabla~\ref{tab:cltv-stats} resume las cinco distribuciones implicadas.
Los valores se han leído directamente de \texttt{cltv.parquet} con
\texttt{pandas.DataFrame.describe} y redondeado a unidades manejables.

\begin{table}[H]
\centering
\caption{Estadísticos por componente de CLTV ($n=5\,750$ clientes; ventana
$W=72$ meses).}
\label{tab:cltv-stats}
\small
\begin{tabular}{l r r r}
\toprule
\textbf{Magnitud} & \textbf{Mediana} & \textbf{Media} & \textbf{Máximo} \\
\midrule
Ingresos brutos (\texteuro)      & $149{,}90$  & $1\,683{,}26$  & $24\,922{,}00$ \\
Margen (proporción)             & $0{,}400$   & $0{,}364$      & $1{,}000$      \\
Frecuencia (tickets)            & $1$         & $3{,}48$       & $34$           \\
Retención $R_t$ ($\,/72$)       & $0{,}0139$  & $0{,}0436$     & $0{,}4028$     \\
\textbf{CLTV (\texteuro)}        & \textbf{0{,}72} & \textbf{$3\,165{,}48$} &
                                   \textbf{$109\,734{,}95$} \\
\bottomrule
\end{tabular}
\end{table}

\noindent La distancia entre mediana y media del CLTV ($0{,}72$\,\texteuro\
frente a $3\,165$\,\texteuro) refleja la asimetría extrema de la
distribución: la mayoría de clientes presenta frecuencia $1$ y retención
$\approx 1/72$, mientras que un pequeño grupo recurrente concentra todo
el valor (figura~\ref{fig:cltv-hist}). El máximo, $109\,735$\,\texteuro,
corresponde al cliente con $\text{customer\_id}=212$.

\begin{figure}[H]
  \centering
  \includegraphics[width=\linewidth]{cltv_distribution.png}
  \caption[Distribución del CLTV]{%
    Distribución del CLTV en \texttt{cltv.parquet}. Histograma en escala
    $\log_{10}$ a la izquierda; \emph{box-plot} con eje horizontal
    en escala simétrica logarítmica (\texttt{symlog}, $\text{linthresh}=1$)
    a la derecha.}
  \label{fig:cltv-hist}
\end{figure}

\section{Métricas adicionales por cliente}

Además del CLTV, el notebook~04 genera tres métricas complementarias en
\texttt{customer\_metrics.parquet}:

\begin{itemize}
  \item \textbf{Recencia} (\texttt{recency\_days}): días desde la última
        compra hasta la fecha de corte. Mediana $421$\,d, P75
        $\approx 1\,300$\,d, P95 $\approx 1\,782$\,d, máximo $1\,824$\,d
        (~$5$ años).
  \item \textbf{Ticket medio} o \emph{AOV} (\texttt{aov}): cociente entre
        ingresos netos y frecuencia. Mediana $129{,}99$\,\texteuro, media
        $210{,}04$\,\texteuro, P75 $\approx 300$\,\texteuro\ y máximo
        $1\,199{,}97$\,\texteuro\ (un AOV negativo aparece para los
        clientes con devolución superior a la venta).
  \item \textbf{Tasa de devolución} (\texttt{return\_rate}):
        $\texttt{returned\_value} / \texttt{gross\_revenue}$. Mediana $0$,
        media $0{,}069$, máximo $2{,}0$ (clientes en los que la cantidad
        devuelta duplica la cantidad vendida; este caso \emph{anómalo} es
        un artefacto del dataset sintético, véase
        capítulo~\ref{ch:conclusiones}).
\end{itemize}

\noindent La tabla~\ref{tab:metricas-extra} ofrece la lectura conjunta de
los estadísticos.

\begin{table}[H]
\centering
\caption{Estadísticos de las métricas adicionales por cliente
($n=5\,750$).}
\label{tab:metricas-extra}
\small
\begin{tabular}{l r r r r}
\toprule
\textbf{Métrica} & \textbf{Mediana} & \textbf{Media} & \textbf{P75} & \textbf{Máx.}\\
\midrule
Recencia (días)         & $421$    & $673{,}5$  & $1\,300$  & $1\,824$    \\
AOV (\texteuro)          & $129{,}99$ & $210{,}04$ & $299{,}99$ & $1\,199{,}97$ \\
Tasa de devolución     & $0$      & $0{,}069$  & $0$       & $2{,}0$     \\
\bottomrule
\end{tabular}
\end{table}

Los tests \texttt{test\_metrics\_aov\_matches\_revenue\_div\_frequency} y
\texttt{test\_recency\_is\_non\_negative} (apéndice~\ref{ap:tests})
formalizan las dos invariantes clave de estas métricas: el AOV recompone
los ingresos netos a partir de la frecuencia, y la recencia nunca puede
ser negativa.
